\chapter{Differentialregning}
\label{ch:diff}
Differentialregning beskriver lokal ændring. Den centrale idé er, at en tilstrækkeligt regulær funktion tæt på et punkt kan erstattes af en affin funktion med en kontrolleret fejl. Hældningen er samtidig en grænseværdi, en lineær afbildning og en lokal følsomhed. Disse tre synsvinkler forbinder de sædvanlige regneregler med fysikkens hastigheder, responskoefficienter og perturbative beskrivelser.

Kapitlet behandler funktioner af én reel variabel. Partielle afledte, Jacobimatricer og differentiabilitet i flere dimensioner hører til kapitel~\ref{ch:flervariabel}. De nævnes kun for at afgrænse, hvilke konklusioner der kan overføres.

\section{Grænseværdier, kontinuitet og lokal adfærd}
\label{sec:graense}
Lad $f$ være defineret i en punkteret omegn af $x_0$. Udsagnet $\lim_{x\to x_0}f(x)=L$ betyder, at der for hvert $\varepsilon>0$ findes et $\delta>0$, sådan at
\[
0<|x-x_0|<\delta\quad\Longrightarrow\quad |f(x)-L|<\varepsilon.
\]
Definitionen kræver ikke, at $f(x_0)$ findes. Kontinuitet i $x_0$ kræver derimod både en funktionsværdi og sammenfaldet $\lim_{x\to x_0}f(x)=f(x_0)$. En venstre- og en højregrænseværdi skal eksistere og være ens, før en endelig tosidet grænseværdi findes. Ved endepunkter arbejder man relativt til definitionsmængden.

Grænseværdier kan føres gennem summer og produkter, når de enkelte endelige grænseværdier findes. For kvotienter skal nævnerens grænseværdi desuden være forskellig fra nul. Ubestemte former som $0/0$ og $\infty-\infty$ er signaler om, at udtrykkets struktur skal undersøges; de er ikke regneregler.

For $h\to0$ betyder $r(h)=o(h)$, at $r(h)/h\to0$. Udsagnet $r(h)=O(h^2)$ betyder, at der findes $C,\delta>0$ med $|r(h)|\le C|h|^2$ for $0<|h|<\delta$. Lille-$o$ udtrykker en forsvindende relativ fejl, mens store-$O$ giver en størrelsesorden med en lokal konstant. Notationen siger ikke noget om en numerisk fejlgrænse, før konstanten er bestemt.

\section{Den afledte og differentialet}
\label{sec:def}
\begin{definition}
Lad $x_0$ være et indre punkt i intervallet $I$. Funktionen $f\colon I\to\R$ er differentiabel i $x_0$, hvis den endelige grænseværdi
\[
 f'(x_0)=\lim_{h\to0}\frac{f(x_0+h)-f(x_0)}{h}
\]
eksisterer. En uendelig hældning tæller ikke som differentiabilitet i denne definition.
\end{definition}
Sekanthældningen måler den gennemsnitlige ændring over et endeligt interval. Den afledte er grænsen af disse hældninger, når intervallet kollapser mod punktet. Den geometriske tangent har ligningen
\[
 y=f(x_0)+f'(x_0)(x-x_0).
\]
Den samme definition kan omformes til
\begin{equation}
 f(x_0+h)=f(x_0)+f'(x_0)h+r(h),\qquad r(h)=o(h).
 \label{eq:linear}
\end{equation}
Denne formulering er særlig vigtig: differentiabilitet er eksistensen af en lokal lineær respons med en fejl, der forsvinder hurtigere end tilvæksten. Differentialet er afbildningen $\dd f_{x_0}\colon h\mapsto f'(x_0)h$. Skrivemåden $\dd f=f'(x)\dd x$ udtrykker denne lineære sammenhæng; den faktiske ændring $\Delta f$ indeholder desuden restleddet.

Hvis $x$ har enhed $[x]$ og $f$ har enhed $[f]$, har $f'$ enhed $[f]/[x]$. Differentialet har samme enhed som $f$. En logaritme eller en eksponentialfunktion skal i en fysisk model have et dimensionsløst argument.

\begin{figure}[htbp]
\centering
\begin{tikzpicture}
\begin{axis}[width=.88\textwidth,height=65mm,axis lines=middle,xlabel={$x$},ylabel={$y$},xmin=-.2,xmax=2.3,ymin=-.8,ymax=4.8,legend pos=north west,legend style={font=\small},samples=100]
\addplot[arkivblaa,thick,domain=0:2.15]{x^2};\addlegendentry{$f(x)=x^2$}
\addplot[arkivrod,thick,domain=0:2.2]{2*x-1};\addlegendentry{Tangent ved $x_0=1$}
\addplot[black,dashed,domain=.4:2.15]{3*x-2};\addlegendentry{Sekant, $h=1$}
\fill[black] (axis cs:1,1) circle[radius=1.5pt];
\fill[black] (axis cs:2,4) circle[radius=1.5pt];
\end{axis}
\end{tikzpicture}
\caption{Sekanten gennem $(1,1)$ og $(2,4)$ har hældning 3; tangenten ved $(1,1)$ har hældning 2. For $h\to0$ går sekantens hældning $2+h$ mod 2. Akserne er dimensionsløse.}
\label{fig:sekant}
\end{figure}


\subsection{Kontinuitet er nødvendig, men ikke tilstrækkelig}
Af \eqref{eq:linear} følger $f(x_0+h)-f(x_0)\to0$, så differentiabilitet medfører kontinuitet. Den omvendte implikation gælder ikke. Funktionen $|x|$ er kontinuert i nul, men har ensidige afledte $-1$ og $1$ og dermed ingen tosidet afledt der.

En afledt behøver heller ikke være kontinuert. Funktionen $f(x)=x^2\sin(1/x)$ for $x\ne0$ og $f(0)=0$ har $f'(0)=0$, mens $f'(x)=2x\sin(1/x)-\cos(1/x)$ for $x\ne0$ ikke har en grænseværdi ved nul. Man må derfor ikke erstatte antagelsen $f\in C^1$ med blot ``$f$ differentiabel'', hvis et argument kræver kontinuitet af $f'$.

\section{Regneregler og sammensatte afhængigheder}
\label{sec:regler}
Nedenstående regler gælder i et punkt, hvor de nævnte funktioner er differentiable. Konstanterne $a,b$ er uafhængige af $x$. For kædereglen skal den indre funktion være differentiabel i $x$, og den ydre funktion i den tilsvarende værdi.

\begin{longtable}{@{}p{0.23\textwidth}p{0.40\textwidth}p{0.29\textwidth}@{}}
\caption{Regneregler. Forudsætninger er en del af formlerne.}\label{tab:regler}\\
\toprule Regel & Formel & Ekstra betingelse\\\midrule\endfirsthead
\toprule Regel & Formel & Ekstra betingelse\\\midrule\endhead
Linearitet & $(af+bg)'=af'+bg'$ & $a,b$ konstante\\[2mm]
Produkt & $(fg)'=f'g+fg'$ & Begge faktorer differentiable\\[2mm]
Kvotient & $\displaystyle (f/g)'=\frac{f'g-fg'}{g^2}$ & $g(x)\ne0$\\[3mm]
Kæderegel & $(f\circ g)'(x)=f'(g(x))g'(x)$ & Sammensætning defineret\\[2mm]
Reciprok & $(1/f)'=-f'/f^2$ & $f(x)\ne0$\\[2mm]
Logaritmisk & $(\ln|f|)'=f'/f$ & $f(x)\ne0$\\[2mm]
Potens af funktion & $(f^\alpha)'=\alpha f^{\alpha-1}f'$ & For reel $\alpha$ er $f>0$ en tilstrækkelig betingelse\\[2mm]
Variabel eksponent & $(f^g)'=f^g(g'\ln f+gf'/f)$ & $f>0$, begge differentiable\\
\bottomrule
\end{longtable}
Produktreglen kommer fra at tilføje og trække $f(x+h)g(x)$ i produktets differens. Efter division med $h$ giver kontinuiteten af $f$ den ene faktor og differentiabiliteten af $g$ den anden. Kædereglen udtrykker, at to lokale lineariseringer sammensættes: den ydre respons skal vurderes ved den indre funktions aktuelle værdi.

Ved tidsafhængighed $q=q(t)$ er $\dd f(q(t))/\dd t=f'(q(t))\dot q(t)$. Prikken betyder differentiation med hensyn til tid, ikke med hensyn til $q$. Denne skelnen forebygger, at en partiel afhængighed forveksles med en total tidsafledt i mekanik og termodynamik.

\subsection{Inverse og implicitte funktioner}
Hvis $f$ er $C^1$ nær $x_0$ og $f'(x_0)\ne0$, er $f$ lokalt invertibel, og
\[
 (f^{-1})'(y_0)=\frac{1}{f'(x_0)},\qquad y_0=f(x_0).
\]
Resultatet er lokalt. En funktion kan godt have flere adskilte monotone grene. Et punkt med $f'=0$ er ikke automatisk et punkt uden en invers, men den inverse kan miste differentiabilitet der.

For en relation $F(x,y)=0$ med $F\in C^1$ og $F_y(x_0,y_0)\ne0$ findes lokalt en funktion $y(x)$, og
\[
 \frac{\dd y}{\dd x}=-\frac{F_x}{F_y}.
\]
Her betegner $F_x,F_y$ partielle afledte, vurderet langs løsningen. Hvis $F_y=0$, er formlen ikke anvendelig; man må undersøge en anden lokal parametrisering eller en mulig singularitet.

\section{Standardafledte med definitionsområder}
\label{sec:standard}
Trigonometriske formler forudsætter radianmål. Ved gradmål kommer en faktor $\pi/180$ ind via kædereglen. De inverse trigonometriske funktioner forstås som deres sædvanlige reelle hovedgrene.

\begin{longtable}{@{}p{0.28\textwidth}p{0.32\textwidth}p{0.30\textwidth}@{}}
\caption{Elementære standardafledte på åbne reelle definitionsområder.}\\
\toprule $f(x)$ & $f'(x)$ & Område eller betingelse\\\midrule\endfirsthead
\toprule $f(x)$ & $f'(x)$ & Område eller betingelse\\\midrule\endhead
$c$ & $0$ & $c$ konstant\\[1mm]
$x^n$ & $nx^{n-1}$ & Heltal $n\ge1$; også nul, hvor udtrykket er defineret\\[1mm]
$x^\alpha$ & $\alpha x^{\alpha-1}$ & $x>0$, $\alpha\in\R$\\[1mm]
$\e^x$ & $\e^x$ & $x\in\R$\\[1mm]
$a^x$ & $a^x\ln a$ & $a>0$ konstant\\[1mm]
$\ln x$ & $1/x$ & $x>0$\\[1mm]
$\ln|x|$ & $1/x$ & $x\ne0$\\[1mm]
$\log_a x$ & $1/(x\ln a)$ & $x>0$, $a>0$, $a\ne1$\\[1mm]
$\sin x$ & $\cos x$ & $x\in\R$\\[1mm]
$\cos x$ & $-\sin x$ & $x\in\R$\\[1mm]
$\tan x$ & $1/\cos^2 x$ & $\cos x\ne0$\\[1mm]
$\cot x$ & $-1/\sin^2 x$ & $\sin x\ne0$\\[1mm]
$\arcsin x$ & $1/\sqrt{1-x^2}$ & $|x|<1$\\[1mm]
$\arccos x$ & $-1/\sqrt{1-x^2}$ & $|x|<1$\\[1mm]
$\arctan x$ & $1/(1+x^2)$ & $x\in\R$\\[1mm]
$\sinh x$ & $\cosh x$ & $x\in\R$\\[1mm]
$\cosh x$ & $\sinh x$ & $x\in\R$\\[1mm]
$\tanh x$ & $1/\cosh^2 x$ & $x\in\R$\\[1mm]
$\operatorname{arsinh}x$ & $1/\sqrt{1+x^2}$ & $x\in\R$\\[1mm]
$\operatorname{arcosh}x$ & $1/\sqrt{x^2-1}$ & $x>1$\\[1mm]
$\operatorname{artanh}x$ & $1/(1-x^2)$ & $|x|<1$\\
\bottomrule
\end{longtable}
En formel gælder kun, hvor den oprindelige funktion er defineret, og de anførte differentiabilitetskrav er opfyldt. Algebraisk forkortelse kan give et udtryk med et større definitionsområde end den oprindelige funktion; det udvider ikke automatisk denne funktion. For $|x|$ er den afledte $\operatorname{sgn}(x)$ for $x\ne0$, men ikke ved nul.

\section{Middelværdisætninger og globale konklusioner}
\label{sec:middel}
\begin{saetning}[Rolle]
Er $f$ kontinuert på $[a,b]$, differentiabel på $(a,b)$ og $f(a)=f(b)$, hvor $a<b$, findes et $c\in(a,b)$ med $f'(c)=0$.
\end{saetning}
Kontinuiteten på det kompakte interval sikrer, at maksimum og minimum antages. Hvis funktionen ikke er konstant, ligger mindst ét af disse ekstrema i det indre. Differentiabilitet tvinger da den afledte til at være nul der. Endepunktskontinuiteten og den indre differentiabilitet udfører altså forskellige opgaver i argumentet.

\begin{saetning}[Lagranges middelværdisætning]
Er $f$ kontinuert på $[a,b]$ og differentiabel på $(a,b)$, findes et $c\in(a,b)$, sådan at
\[
 f(b)-f(a)=f'(c)(b-a).
\]
\end{saetning}
Sætningen følger ved at trække sekantens affine funktion fra $f$ og anvende Rolle. Den siger, at den gennemsnitlige ændring svarer til en øjeblikkelig ændring et sted i intervallet. Den angiver generelt hverken et entydigt $c$ eller en konstruktiv metode til at finde det.

Hvis $|f'(x)|\le M$ på intervallet, følger stabilitetsestimatet
\[
 |f(b)-f(a)|\le M|b-a|.
\]
Det er en Lipschitzgrænse: en begrænset lokal følsomhed giver en kontrolleret endelig ændring. Hvis $f'\ge0$ på et interval, er $f$ ikke-aftagende. Hvis $f'>0$ overalt, er $f$ strengt voksende. Den omvendte slutning er stærkere end nødvendigt: en strengt voksende differentiabel funktion kan have $f'=0$ i enkelte punkter.

Cauchys middelværdisætning gælder for to funktioner med de samme kontinuitets- og differentiabilitetskrav:
\[
 [f(b)-f(a)]g'(c)=[g(b)-g(a)]f'(c)
\]
for et $c\in(a,b)$. Man må først dividere med $g'(c)$ og $g(b)-g(a)$, når de er forskellige fra nul. Denne sætning ligger bag l'Hôpitals regel.

\section{Højere afledte, monotoni og krumning}
\label{sec:hoejere}
Den anden afledte $f''$ beskriver ændringen i hældningen. Klassen $C^n(I)$ består af funktioner med kontinuerte afledte til og med orden $n$ på $I$. $C^\infty$ betyder glat, men betyder ikke nødvendigvis analytisk: en funktion kan være uendeligt differentiabel uden at være lig sin Taylorrække.

\subsection{Ekstrema og stationære punkter}
For en differentiabel funktion er $f'(x_0)=0$ en nødvendig betingelse for et indre lokalt ekstremum. Den er ikke tilstrækkelig. Fortegnsskiftet af $f'$ afgør typen: fra negativ til positiv giver et minimum; fra positiv til negativ giver et maksimum. Intet fortegnsskift giver normalt intet ekstremum, men degenererede tilfælde kræver en direkte undersøgelse.

Hvis $f\in C^2$ nær $x_0$, $f'(x_0)=0$ og $f''(x_0)>0$, er punktet et strengt lokalt minimum; ved $f''(x_0)<0$ et strengt lokalt maksimum. Hvis $f''(x_0)=0$, er testen uafgjort. Mere generelt: er den første ikke-nul afledte efter konstantleddet af orden $m$, og gælder en Taylorudvikling med restled $o((x-x_0)^m)$, giver lige $m$ et ekstremum bestemt af fortegnet; ulige $m$ giver intet ekstremum.

En fuldstændig undersøgelse på et lukket interval skal også omfatte endepunkter og indre punkter, hvor funktionen ikke er differentiabel. En liste over løsninger til $f'=0$ alene er ikke en optimeringsmetode på hele definitionsmængden.

\subsection{Konveksitet og geometrisk krumning}
En differentiabel funktion er konveks på et interval, når dens graf ligger over alle sine tangenter. For $C^2$-funktioner er dette ækvivalent med $f''\ge0$. Konkavitet svarer til $f''\le0$. Et vendepunkt kræver skift mellem konveksitet og konkavitet; ligningen $f''=0$ alene er utilstrækkelig.

For en plan graf $y=f(x)$ med dimensionsløse, lige skalerede koordinater er den geometriske krumnings størrelse
\[
 \kappa(x)=\frac{|f''(x)|}{(1+f'(x)^2)^{3/2}}.
\]
Andenafledte og geometrisk krumning er altså forskellige størrelser. Hvis akserne har forskellige fysiske enheder, skal en fysisk indlejring og skalering fastlægges, før denne geometriske formel kan fortolkes.

\begin{figure}[htbp]
\centering
\begin{tikzpicture}
\begin{axis}[width=.88\textwidth,height=65mm,axis lines=middle,xlabel={$x$},ylabel={$y$},xmin=-2,xmax=2,ymin=-3.5,ymax=3.5,legend pos=north west,legend style={font=\small},samples=150]
\addplot[arkivblaa,thick,domain=-1.9:1.9]{x^3-3*x};\addlegendentry{$f=x^3-3x$}
\addplot[arkivrod,dashed,domain=-1.45:1.45]{3*x^2-3};\addlegendentry{$f'=3x^2-3$}
\fill[black] (axis cs:-1,2) circle[radius=1.5pt];
\fill[black] (axis cs:1,-2) circle[radius=1.5pt];
\fill[black] (axis cs:0,0) circle[radius=1.5pt];
\end{axis}
\end{tikzpicture}
\caption{Afledtens fortegn styrer monotonien. Ved $x=-1$ ligger et maksimum, ved $x=1$ et minimum. Vendepunktet $x=0$ har $f''=0$ med fortegnsskift, men er ikke stationært. Akserne er dimensionsløse.}
\end{figure}


\section{Taylorudvikling og kontrollerede restled}
\label{sec:taylor}
Taylorpolynomiet af grad $n$ omkring $a$ er
\[
 T_n(x)=\sum_{k=0}^n\frac{f^{(k)}(a)}{k!}(x-a)^k.
\]
Det matcher funktionsværdien og de første $n$ afledte i udviklingspunktet. Polynomiet beskriver lokal struktur; det er ikke en global erstatning uden en fejlundersøgelse.

\begin{saetning}[Taylor med Lagrangerest]
Lad $f\in C^{n+1}$ på et åbent interval, der indeholder segmentet mellem $a$ og $x$. Da findes et $\xi$ mellem $a$ og $x$, sådan at
\begin{equation}
 f(x)=T_n(x)+\frac{f^{(n+1)}(\xi)}{(n+1)!}(x-a)^{n+1}.
 \label{eq:rest}
\end{equation}
\end{saetning}
Den valgte $C^{n+1}$-antagelse er en enkel tilstrækkelig betingelse. Hvis $|f^{(n+1)}(t)|\le M$ på segmentet, fås den praktiske fejlgrænse
\[
 |f(x)-T_n(x)|\le\frac{M|x-a|^{n+1}}{(n+1)!}.
\]
Ved blot $f\in C^n$ nær $a$ findes Peanoformen $f(x)=T_n(x)+o((x-a)^n)$ for $x\to a$. Den form giver ikke i sig selv en kendt konstant for fejlens størrelse. Når $f^{(n+1)}$ er kontinuert, er også integralrestleddet gyldigt:
\[
 R_n(x)=\frac{1}{n!}\int_a^x (x-t)^n f^{(n+1)}(t)\dd t.
\]
Integralformen er velegnet til fortegns- og normestimater og gælder også for $x<a$ med integralets sædvanlige orientering.

\subsection{Standardudviklinger omkring nul}
Her er $z$ dimensionsløs, og restleddene gælder for $z\to0$:
\begin{align*}
 \e^z&=1+z+\tfrac12z^2+\tfrac16z^3+O(z^4),\\
 \sin z&=z-\tfrac16z^3+O(z^5),\\
 \cos z&=1-\tfrac12z^2+\tfrac1{24}z^4+O(z^6),\\
 \ln(1+z)&=z-\tfrac12z^2+\tfrac13z^3+O(z^4),\\
 (1+z)^\alpha&=1+\alpha z+\tfrac12\alpha(\alpha-1)z^2+O(z^3).
\end{align*}
Logaritmen og den reelle potens kræver $1+z>0$. De tilhørende uendelige potensrækker har yderligere konvergensbetingelser: logaritme- og den generelle binomialrække har konvergensradius 1, med særskilt analyse på randen; for ikke-negative heltalseksponenter stopper binomialrækken og er et polynomium. Eksponential-, sinus- og cosinusrækkerne konvergerer for alle reelle $z$.

At alle afledte findes, sikrer ikke, at en uendelig Taylorrække repræsenterer funktionen. Funktionen $\e^{-1/x^2}$ for $x\ne0$, udvidet med værdien nul i nul, er glat og har alle afledte lig nul i nul, men er positiv uden for nul. Analyticitet er en ekstra egenskab.

\begin{figure}[htbp]
\centering
\begin{tikzpicture}
\begin{axis}[width=.88\textwidth,height=65mm,axis lines=middle,xlabel={$x$},ylabel={$y$},xmin=-2,xmax=2,ymin=-1.2,ymax=1.2,legend pos=north west,legend style={font=\small},samples=150]
\addplot[arkivblaa,thick,domain=-2:2]{sin(deg(x))};\addlegendentry{$\sin x$}
\addplot[black,dotted,domain=-1.2:1.2]{x};\addlegendentry{$T_1=x$}
\addplot[arkivrod,dashed,thick,domain=-2:2]{x-x^3/6};\addlegendentry{$T_3=x-x^3/6$}
\end{axis}
\end{tikzpicture}
\caption{Taylorpolynomier omkring nul. Den kubiske tilnærmelse følger sinus tættere nær nul end tangenten. For alle reelle $x$ giver Taylor til orden 4, hvor koefficienten til $x^4$ er nul, grænsen $|\sin x-T_3(x)|\le |x|^5/120$. En lille relativ fejl følger ikke automatisk nær sinusfunktionens nulpunkter.}
\end{figure}


\section{Grænseværdier og l'Hôpitals regel}
\label{sec:hopital}
En anvendelig version er følgende: Lad $f,g$ være differentiable på et punkteret ensidigt interval mod et endepunkt $a$, med $g'\ne0$. Antag enten $f,g\to0$ eller $|f|,|g|\to\infty$. Hvis $f'/g'$ har en grænseværdi $L$, endelig eller uendelig, gælder under disse betingelser
\[
 \lim_{x\to a}\frac{f(x)}{g(x)}=L.
\]
Samme princip kan anvendes ved uendelige endepunkter. Tosidige grænseværdier kræver kontrol fra begge sider. Ved gentagen brug skal betingelserne verificeres igen på hvert trin. Hvis kvotienten af de afledte ikke har en grænseværdi, er testen uafgjort; det beviser ikke, at den oprindelige kvotient mangler en grænseværdi.

Reglen differentierer tæller og nævner hver for sig til en ny kvotient; den beregner ikke den afledte af den oprindelige kvotient. Former som $0\cdot\infty$ skal først omskrives til en relevant kvotient. Taylorudvikling er ofte mere informativ, fordi den samtidig viser den første ikke-nul orden og tilnærmelsens fejl.

\section{Fysisk fortolkning og følsomhed}
\label{sec:fysik}
For en tilstrækkeligt glat position $q(t)$ er hastigheden $v=\dot q$ og accelerationen $a=\ddot q$. De har enheder henholdsvis længde pr. tid og længde pr. tid i anden. En grafisk hældning afhænger af aksernes skala, mens den fysiske afledte bestemmes af variablerne og deres enheder.

For en éndimensional konservativ kraft er $F(q)=-U'(q)$. Et ligevægtspunkt opfylder $U'(q_0)=0$. Hvis $U''(q_0)>0$ og massen $m>0$, giver lineariseringen
\[
 m\ddot\eta=-U''(q_0)\eta+O(\eta^2),\qquad \eta=q-q_0,
\]
forudsat tilstrækkelig regularitet, eksempelvis $U\in C^3$ nær $q_0$. Den lineære svingningsfrekvens er $\omega=\sqrt{U''(q_0)/m}$. For $U''(q_0)<0$ er den lineariserede ligevægt ustabil. Ved $U''(q_0)=0$ må højere ordener undersøges. Dette er en generel kobling mellem afledte og stabilitet, ikke en konkret opgaveløsning.

For en målt variabel $x$ giver en lille perturbation $\delta x$ den førsteordensændring $\delta f\approx f'(x)\delta x$. Hvis $x\ne0$ og $f(x)\ne0$, er den relative følsomhed
\[
 \kappa_{\mathrm{rel}}(x)=\left|\frac{x f'(x)}{f(x)}\right|,\qquad
 \left|\frac{\delta f}{f}\right|\approx\kappa_{\mathrm{rel}}\left|\frac{\delta x}{x}\right|.
\]
Nær nulpunkter kan relativ fejl være uegnet, selv når den absolutte fejl er lille. For en stokastisk usikkerhed med lille standardafvigelse $\sigma_x$ giver lineariseringen $\sigma_f\approx|f'(x)|\sigma_x$. Ikke-linearitet, skæve fordelinger og korrelationer kræver en videre analyse i det selvstændige statistikkompendium; formlen er ikke en generel eksakt usikkerhedslov.

\section{Numerisk differentiation og fejlbalance}
\label{sec:numerisk}
Når $f$ kun er tilgængelig gennem evalueringer, kan differenskvotienter approksimere den afledte. For $h\ne0$ er de mest brugte:
\begin{align*}
 D_h^+f(x)&=\frac{f(x+h)-f(x)}{h}
           =f'(x)+\tfrac h2f''(x)+O(h^2),\\
 D_h^0f(x)&=\frac{f(x+h)-f(x-h)}{2h}
           =f'(x)+\tfrac{h^2}{6}f^{(3)}(x)+O(h^4),\\
 D_h^2f(x)&=\frac{f(x+h)-2f(x)+f(x-h)}{h^2}
           =f''(x)+\tfrac{h^2}{12}f^{(4)}(x)+O(h^4).
\end{align*}
Tilstrækkelige lokale regularitetskrav for de viste restled er henholdsvis $C^3$, $C^5$ og $C^6$. De centrerede formler udnytter symmetri til at eliminere ulige eller lige bidrag. Ved rande kræver de funktionsværdier uden for intervallet; brug i stedet en passende ensidig formel eller en fysisk begrundet randudvidelse.

Mindre $h$ er ikke altid bedre. Hvis hver funktionsværdi har absolut fejl højst $\varepsilon_f$, kan fejlbidraget til den centrerede førsteafledte være op til $\varepsilon_f/|h|$. En simpel model for totalfejlen er derfor
\[
 E(h)\lesssim C h^2+\frac{\varepsilon_f}{|h|}.
\]
Her kommer første led fra trunkering og andet led fra afrunding eller datastøj. Balancen giver skalaen $|h|\sim(\varepsilon_f/C)^{1/3}$ op til konstante faktorer, når $C>0$. Dette er et lokalt skalaestimat, ikke et universelt optimalt skridt. Variabelskalering, støjstruktur og ukendte højere afledte påvirker valget.

\begin{figure}[htbp]
\centering
\begin{tikzpicture}
\begin{loglogaxis}[width=.88\textwidth,height=65mm,xlabel={Skridt $h$},ylabel={Fejlmodel},xmin=1e-7,xmax=1e-1,ymin=1e-11,ymax=1e-1,legend pos=north east,legend style={font=\small},samples=100]
\addplot[arkivblaa,dashed,domain=1e-7:1e-1]{x^2};\addlegendentry{$h^2$}
\addplot[arkivrod,dotted,thick,domain=1e-7:1e-1]{1e-12/x};\addlegendentry{$10^{-12}/h$}
\addplot[black,thick,domain=1e-7:1e-1]{x^2+1e-12/x};\addlegendentry{Sum}
\end{loglogaxis}
\end{tikzpicture}
\caption{Illustrativ fejlmodel for centreret differentiation med $C=1$ og $\varepsilon_f=10^{-12}$ i dimensionsløse variable. Minimum ligger ved $h=(\varepsilon_f/2C)^{1/3}\approx7{,}94\cdot10^{-5}$. Kurverne er en analytisk model, ikke målte flydende-kommafejl.}
\end{figure}


Automatisk differentiation beregner afledte ved at føre kædereglen gennem programmets operationer. Den undgår differenskvotientens trunkeringsfejl, men hverken diskontinuerte programgrene, afrunding eller en forkert fysisk model forsvinder af den grund. Symbolsk differentiation, automatisk differentiation og numeriske differenser er forskellige redskaber.

\clearpage
\section{Metodeoverblik og kontrolpunkter}
\label{sec:kontrol}
Begynd med definitionsmængde og ønsket konklusion. Til lokal linearisering bruges differentiabilitet; til en kvantitativ fejlgrænse kræves kontrol over relevante højere afledte. Til global monotoni anvendes fortegnet på den afledte på hele det relevante interval. Til ekstremumskontrol undersøges både stationære punkter, singulariteter og rande.

Ved sammensætninger skal den indre funktions afledte og den ydre funktions evalueringspunkt fremgå. Ved inverse funktioner angives gren og ikke-nul hældning. Ved Taylorapproksimation angives udviklingspunkt, orden, gyldighedsområde og restled. Ved numerisk differentiation afprøves flere skridtstørrelser, så både konvergensområdet og støjgulvet kan identificeres.

En faglig kontrol omfatter desuden:
\begin{itemize}
\item Har alle symboler en definition, og er variable og konstanter skilt ad?
\item Stemmer enhederne, også efter tidsdifferentiation og kædereglen?
\item Er udsagnet lokalt eller globalt, og ligger punktet i det indre eller på randen?
\item Er der divideret med en størrelse, som kan være nul?
\item Er en nødvendig betingelse fejlagtigt brugt som tilstrækkelig?
\item Matcher figurernes akser, kurver og markeringer de beskrevne formler?
\end{itemize}
Disse kontroller er en del af den matematiske argumentation. De erstatter ikke et bevis, men gør skjulte antagelser synlige og forbinder håndbogens formler med forsvarlig anvendelse i fysik.
